\section{Related Work}
\label{sec:related}

\textbf{Hallucination.} Surveys define hallucination as generated content that is fluent but unfaithful to the source or to world knowledge~\cite{ji2023survey,huang2025survey}. Benchmarks such as HaluEval~\cite{li2023halueval} and FActScore~\cite{min2023factscore} measure it in dialogue, summarisation and long-form biographies. Theoretical work shows that it cannot be eliminated: a calibrated generative model must hallucinate on facts that appear rarely in its training data~\cite{kalai2024calibrated}, and no computable model can avoid it on all inputs~\cite{xu2024inevitable}. Kalai \emph{et al.}~\cite{kalai2025why} further argue that accuracy-only training and evaluation reward guessing over admitting uncertainty, and fine-tuning on facts that a model did not already know increases its tendency to hallucinate~\cite{gekhman2024finetuning}. Together these results shift the practical question from \emph{whether} a model hallucinates to whether it can tell \emph{when} it is likely to.

\textbf{Unanswerable and false-premise questions.} SQuAD~2.0 introduced unanswerable questions for extractive reading comprehension~\cite{rajpurkar2018squad2}. In open-domain settings, (QA)$^2$~\cite{kim2023qa2}, CREPE~\cite{yu2023crepe} and FalseQA~\cite{hu2023falseqa} collect questions with questionable or false presuppositions, and TruthfulQA probes imitative falsehoods~\cite{lin2022truthfulqa}. SelfAware~\cite{yin2023selfaware} and KUQ~\cite{amayuelas2024kuq} test whether models recognise unanswerable questions. AbstentionBench~\cite{kirichenko2025abstentionbench} finds that abstention remains unsolved across 20 frontier models and does not improve reliably with scale, and CoCoNot~\cite{brahman2024noncompliance} broadens the problem to contextual noncompliance. These benchmarks pool heterogeneous questions, so answerable and unanswerable items differ in topic, length and style as well as answerability.

\textbf{Knowledge boundaries and abstention.} Accuracy drops sharply for less popular entities~\cite{mallen2023popqa} and for facts that are seldom mentioned in pre-training text~\cite{kandpal2023longtail}, and LLMs remain overconfident about such facts~\cite{ren2025boundary}; see~\cite{li2025boundary} for a survey. Selective question answering abstains when a calibrator predicts an error~\cite{kamath2020selective}. More recent work fine-tunes models to say ``I don't know'', using refusal-aware tuning~\cite{zhang2024rtuning}, model-specific ``Idk'' data~\cite{cheng2024idk} or honesty-oriented alignment~\cite{yang2024honesty}; the last of these explicitly targets refusing unknown questions \emph{without} becoming over-conservative. Other work detects knowledge gaps through multi-model collaboration~\cite{feng2024abstain}; see~\cite{wen2025abstention} for a survey.

\textbf{Uncertainty estimation and calibration.} Calibration is classically measured with the expected calibration error (ECE)~\cite{naeini2015ece,guo2017calibration}, and early studies found that pre-trained QA models are poorly calibrated~\cite{jiang2021calibration}. For LLMs, uncertainty can be read from token probabilities, from self-evaluation such as \ptrue{} and \pik{}~\cite{kadavath2022know}, or from verbalised confidence~\cite{lin2022words,mielke2022linguistic,tian2023just,xiong2024express}. Verbal expressions of certainty may, however, imitate how humans write rather than reflect the model's own uncertainty~\cite{zhou2023grey}, and prompting alone yields poorly calibrated estimates~\cite{kapoor2024taught}. A third family uses the agreement of sampled answers, an idea that traces back to self-consistency~\cite{wang2023selfconsistency}. Examples include SelfCheckGPT~\cite{manakul2023selfcheckgpt}, semantic entropy~\cite{kuhn2023semantic,farquhar2024nature}, repeated sampling for ambiguous questions~\cite{cole2023selectively}, and combinations of consistency with self-reflection~\cite{chen2024bsdetector}. A related line shows that models ``know more than they say''. Hidden states encode truthfulness~\cite{azaria2023lying,orgad2025know} and even answerability~\cite{slobodkin2023curious}, and models can recognise many of their own errors when asked about them separately~\cite{zhang2024snowball}. Most of this work evaluates models of 7B parameters or more and uses answerable questions or questions with a supporting context.

\textbf{Research gap.} To our knowledge, no prior study \emph{jointly} (i)~contrasts unanswerable questions with minimally different answerable twins, which controls for topic and surface form; (ii)~places fabricated entities on the same popularity continuum as real long-tail entities; and (iii)~compares what a model \emph{says} (refusals, verbal confidence) with what its probabilities \emph{reveal} (token likelihood, sampling consistency, \ptrue{}, \pik{}), all without training and with access only to output probabilities. It also remains unclear whether these findings hold for small open-weight models (${\le}3$B parameters), which are increasingly deployed offline on commodity hardware. This paper addresses all three points for such models.
